% GENERATED FILE. Edit canonical lessons, then run npm run generate:books.
% canonical-source-hash: fnv1a64:bf0434d634aa301b
% canonical-lessons: ML-C372-tattippu, ML-C372-tattippukaran, ML-C372-vancikkuka, ML-C372-parrikkuka, ML-C372-kabalippikkuka, ML-R372-review-365-372

\chapter{Fraud and Tricks}
\label{ch:ml-fraud-and-tricks}
\hlchaptermodality{372}

\begin{chapteropening}
\textbf{By the end of this chapter:} \emph{I can say a fraud, a swindler, to deceive, to trick and to fool.}

Complete the last lesson of chapter 372: I can say a fraud, a swindler, to deceive, to trick and to fool.
\end{chapteropening}

\section[\texorpdfstring{\ml{തട്ടിപ്പ്} (taṭṭippŭ)}{തട്ടിപ്പ് (taṭṭippŭ)}]{\ml{തട്ടിപ്പ്} (taṭṭippŭ) — a fraud; a scam}

\label{lesson:ML-C372-tattippu}



\begin{quote}
Before the new one: say the Malayalam for to release, then the Malayalam for to free.
\end{quote}

\subsection*{You'll want to know: \ml{തട്ടിപ്പ്}}
\textbf{\ml{തട്ടിപ്പ്}} — \emph{taṭṭippŭ} — \textquotedblleft{}a fraud; a scam\textquotedblright{}.

\begin{quote}
\ml{ഫോണിലൂടെയുള്ള} \ml{തട്ടിപ്പ്} \ml{ഇപ്പോൾ} \ml{കൂടുതലാണ്}.
\end{quote}

\begin{quote}
\emph{phōṇilūṭeyuḷḷa taṭṭippŭ ippōḷ kūṭutalāṇŭ.}
\end{quote}

\begin{quote}
Fraud over the phone is more common now.
\end{quote}

A scam takes your money by a trick.

\begin{cousinweb}
\textbf{\ml{തട്ടിപ്പ്}} (\emph{taṭṭippŭ}) is native Dravidian, a noun built on \ml{തട്ടുക} (\emph{taṭṭuka}), \textquotedblleft{}to knock, to strike\textquotedblright{} (DEDR 3039), which in everyday speech also means to cheat someone of money.
\end{cousinweb}

\begin{grammarlens}[title={What you've built}]
One more word for telling what happened.
\end{grammarlens}

\subsection*{Your turn}
\begin{itemize}
\raggedright
  \item \emph{Say it:} \emph{taṭṭippŭ}
  \item \emph{Say it:} \emph{taṭṭippŭ}, once more
  \item \emph{Say it:} \emph{phōṇilūṭeyuḷḷa taṭṭippŭ ippōḷ kūṭutalāṇŭ.} — Fraud over the phone is more common now.
  \item \emph{Recall:} say the Malayalam for to release, then the Malayalam for to free, then say \emph{taṭṭippŭ} again
\end{itemize}

\begin{tcolorbox}[breakable,colback=teal!4,colframe=teal!35!black,title={Before you move on}]
What does \ml{തട്ടിപ്പ്} mean? (\textquotedblleft{}A fraud; a scam\textquotedblright{}.) Say it once more.
\end{tcolorbox}
\section[\texorpdfstring{\ml{തട്ടിപ്പുകാരൻ} (taṭṭippukāran)}{തട്ടിപ്പുകാരൻ (taṭṭippukāran)}]{\ml{തട്ടിപ്പുകാരൻ} (taṭṭippukāran) — a swindler; a con man}

\label{lesson:ML-C372-tattippukaran}



\begin{quote}
Before the new one: say the Malayalam for to free, then the Malayalam for a fraud.
\end{quote}

\subsection*{You'll want to know: \ml{തട്ടിപ്പുകാരൻ}}
\textbf{\ml{തട്ടിപ്പുകാരൻ}} — \emph{taṭṭippukāran} — \textquotedblleft{}a swindler; a con man\textquotedblright{}.

\begin{quote}
\ml{ആ} \ml{തട്ടിപ്പുകാരൻ} \ml{ഒരുപാട്} \ml{ആളുകളുടെ} \ml{പണം} \ml{കൊണ്ടുപോയി}.
\end{quote}

\begin{quote}
\emph{ā taṭṭippukāran orupāṭŭ āḷukaḷuṭe paṇaṁ koṇṭupōyi.}
\end{quote}

\begin{quote}
That swindler made off with many people's money.
\end{quote}

Plural: \ml{തട്ടിപ്പുകാർ}, \emph{taṭṭippukār}.

\ml{തട്ടിപ്പ്} + -\ml{കാരൻ}, one who does it.

\begin{cousinweb}
\textbf{\ml{തട്ടിപ്പുകാരൻ}} (\emph{taṭṭippukāran}) is native Dravidian: \ml{തട്ടിപ്പ്} (\emph{taṭṭippŭ}), \textquotedblleft{}a fraud\textquotedblright{}, from \ml{തട്ടുക} (DEDR 3039), with the ending -\ml{കാരൻ}, \textquotedblleft{}one who does it\textquotedblright{}.
\end{cousinweb}

\begin{grammarlens}[title={What you've built}]
One more word for telling what happened.
\end{grammarlens}

\subsection*{Your turn}
\begin{itemize}
\raggedright
  \item \emph{Say it:} \emph{taṭṭippukāran}
  \item \emph{Say it:} \emph{taṭṭippukāran}, once more
  \item \emph{Say it:} \emph{ā taṭṭippukāran orupāṭŭ āḷukaḷuṭe paṇaṁ koṇṭupōyi.} — That swindler made off with many people's money.
  \item \emph{Recall:} say the Malayalam for to free, then the Malayalam for a fraud, then say \emph{taṭṭippukāran} again
\end{itemize}

\begin{tcolorbox}[breakable,colback=teal!4,colframe=teal!35!black,title={Before you move on}]
What does \ml{തട്ടിപ്പുകാരൻ} mean? (\textquotedblleft{}A swindler; a con man\textquotedblright{}.) Say it once more.
\end{tcolorbox}
\section[\texorpdfstring{\ml{വഞ്ചിക്കുക} (vañcikkuka)}{വഞ്ചിക്കുക (vañcikkuka)}]{\ml{വഞ്ചിക്കുക} (vañcikkuka) — to deceive; to betray}

\label{lesson:ML-C372-vancikkuka}



\begin{quote}
Before the new one: say the Malayalam for a fraud, then the Malayalam for a swindler.
\end{quote}

\subsection*{You'll want to know: \ml{വഞ്ചിക്കുക}}
\textbf{\ml{വഞ്ചിക്കുക}} — \emph{vañcikkuka} — \textquotedblleft{}to deceive; to betray\textquotedblright{}.

\begin{quote}
\ml{അവൻ} \ml{സ്വന്തം} \ml{സുഹൃത്തിനെ} \ml{വഞ്ചിച്ചു}.
\end{quote}

\begin{quote}
\emph{avan svantaṁ suhṛttine vañciccu.}
\end{quote}

\begin{quote}
He deceived his own friend.
\end{quote}

Past: \ml{വഞ്ചിച്ചു}, \emph{vañciccu}.

The person deceived takes -\ml{എ}: \ml{സുഹൃത്തിനെ} \ml{വഞ്ചിച്ചു}.

\begin{cousinweb}
\textbf{\ml{വഞ്ചിക്കുക}} (\emph{vañcikkuka}) is built from Sanskrit \emph{vañcana}, \textquotedblleft{}cheating, deception, fraud\textquotedblright{} (Monier-Williams), with the Malayalam verb ending -\ml{ിക്കുക}.
\end{cousinweb}

\begin{grammarlens}[title={What you've built}]
One more word for telling what happened.
\end{grammarlens}

\subsection*{Your turn}
\begin{itemize}
\raggedright
  \item \emph{Say it:} \emph{vañcikkuka}
  \item \emph{Say it:} \emph{vañcikkuka}, once more
  \item \emph{Say it:} \emph{avan svantaṁ suhṛttine vañciccu.} — He deceived his own friend.
  \item \emph{Recall:} say the Malayalam for a fraud, then the Malayalam for a swindler, then say \emph{vañcikkuka} again
\end{itemize}

\begin{tcolorbox}[breakable,colback=teal!4,colframe=teal!35!black,title={Before you move on}]
What does \ml{വഞ്ചിക്കുക} mean? (\textquotedblleft{}To deceive; to betray\textquotedblright{}.) Say it once more.
\end{tcolorbox}
\section[\texorpdfstring{\ml{പറ്റിക്കുക} (paṟṟikkuka)}{പറ്റിക്കുക (paṟṟikkuka)}]{\ml{പറ്റിക്കുക} (paṟṟikkuka) — to trick; to con}

\label{lesson:ML-C372-parrikkuka}



\begin{quote}
Before the new one: say the Malayalam for a swindler, then the Malayalam for to deceive.
\end{quote}

\subsection*{You'll want to know: \ml{പറ്റിക്കുക}}
\textbf{\ml{പറ്റിക്കുക}} — \emph{paṟṟikkuka} — \textquotedblleft{}to trick; to con\textquotedblright{}.

\begin{quote}
\ml{ആ} \ml{കച്ചവടക്കാരൻ} \ml{എന്നെ} \ml{പറ്റിച്ചു}.
\end{quote}

\begin{quote}
\emph{ā kaccavaṭakkāran enne paṟṟiccu.}
\end{quote}

\begin{quote}
That trader conned me.
\end{quote}

Past: \ml{പറ്റിച്ചു}, \emph{paṟṟiccu}.

The everyday word; \ml{വഞ്ചിക്കുക} is stronger and more formal.

\begin{cousinweb}
\textbf{\ml{പറ്റിക്കുക}} (\emph{paṟṟikkuka}) is native Dravidian, usually traced to \ml{പറ്റുക} (\emph{paṟṟuka}), \textquotedblleft{}to stick to, to catch\textquotedblright{} (DEDR 4034), the verb inside \ml{പറ്റില്ല}, \textquotedblleft{}it is not possible\textquotedblright{}.
\end{cousinweb}

\begin{grammarlens}[title={What you've built}]
One more word for telling what happened.
\end{grammarlens}

\subsection*{Your turn}
\begin{itemize}
\raggedright
  \item \emph{Say it:} \emph{paṟṟikkuka}
  \item \emph{Say it:} \emph{paṟṟikkuka}, once more
  \item \emph{Say it:} \emph{ā kaccavaṭakkāran enne paṟṟiccu.} — That trader conned me.
  \item \emph{Recall:} say the Malayalam for a swindler, then the Malayalam for to deceive, then say \emph{paṟṟikkuka} again
\end{itemize}

\begin{tcolorbox}[breakable,colback=teal!4,colframe=teal!35!black,title={Before you move on}]
What does \ml{പറ്റിക്കുക} mean? (\textquotedblleft{}To trick; to con\textquotedblright{}.) Say it once more.
\end{tcolorbox}
\section[\texorpdfstring{\ml{കബളിപ്പിക്കുക} (kabaḷippikkuka)}{കബളിപ്പിക്കുക (kabaḷippikkuka)}]{\ml{കബളിപ്പിക്കുക} (kabaḷippikkuka) — to fool; to dupe}

\label{lesson:ML-C372-kabalippikkuka}



\begin{quote}
Before the new one: say the Malayalam for to deceive, then the Malayalam for to trick.
\end{quote}

\subsection*{You'll want to know: \ml{കബളിപ്പിക്കുക}}
\textbf{\ml{കബളിപ്പിക്കുക}} — \emph{kabaḷippikkuka} — \textquotedblleft{}to fool; to dupe\textquotedblright{}.

\begin{quote}
\ml{അവൻ} \ml{വ്യാജ} \ml{സർട്ടിഫിക്കറ്റ്} \ml{കാണിച്ച്} \ml{എല്ലാവരെയും} \ml{കബളിപ്പിച്ചു}.
\end{quote}

\begin{quote}
\emph{avan vyāja sarṭṭiphikkaṟṟŭ kāṇiccŭ ellāvareyuṁ kabaḷippiccu.}
\end{quote}

\begin{quote}
He fooled everyone by showing a fake certificate.
\end{quote}

Past: \ml{കബളിപ്പിച്ചു}, \emph{kabaḷippiccu}.

Formal; common in news reports of fraud.

\begin{cousinweb}
The origin of \textbf{\ml{കബളിപ്പിക്കുക}} (\emph{kabaḷippikkuka}) is not settled. It carries the Malayalam causative ending -\ml{ിപ്പിക്കുക}: to make someone believe what is false.
\end{cousinweb}

\begin{grammarlens}[title={What you've built}]
One more word for telling what happened.
\end{grammarlens}

\subsection*{Your turn}
\begin{itemize}
\raggedright
  \item \emph{Say it:} \emph{kabaḷippikkuka}
  \item \emph{Say it:} \emph{kabaḷippikkuka}, once more
  \item \emph{Say it:} \emph{avan vyāja sarṭṭiphikkaṟṟŭ kāṇiccŭ ellāvareyuṁ kabaḷippiccu.} — He fooled everyone by showing a fake certificate.
  \item \emph{Recall:} say the Malayalam for to deceive, then the Malayalam for to trick, then say \emph{kabaḷippikkuka} again
\end{itemize}

\begin{tcolorbox}[breakable,colback=teal!4,colframe=teal!35!black,title={Before you move on}]
What does \ml{കബളിപ്പിക്കുക} mean? (\textquotedblleft{}To fool; to dupe\textquotedblright{}.) Say it once more.
\end{tcolorbox}
\section[(dialogue)]{Words from chapters 365–372 — a review}

\label{lesson:ML-R372-review-365-372}



\begin{quote}
Say the Malayalam for a complaint and a criminal.
\end{quote}

\subsection*{Your turn}
Say these aloud:

\begin{itemize}
\raggedright
  \item a complaint, a criminal, an investigation, to arrest, to seize
  \item a clue, a mark, a footprint, a fingerprint, mysterious
  \item a crime, a confession, to capture, to kidnap, to keep watch on
  \item a court of law, a court case, a trial, a judge, the accused
  \item a verdict, the guilty one, bail, an innocent person, to pass sentence
  \item testimony, an eyewitness, a lie, to testify, to swear an oath
  \item a jail, imprisonment, a prisoner, to release, to free
  \item a fraud, a swindler, to deceive, to trick, to fool
\end{itemize}

\begin{tcolorbox}[breakable,colback=teal!4,colframe=teal!35!black,title={Before you move on}]
A complaint? (\textbf{\ml{പരാതി}}, \emph{parāti}.) A criminal? (\textbf{\ml{കുറ്റവാളി}}, \emph{kuṟṟavāḷi}.) An investigation? (\textbf{\ml{അന്വേഷണം}}, \emph{anvēṣaṇaṁ}.) To arrest? (\textbf{\ml{അറസ്റ്റ്} \ml{ചെയ്യുക}}, \emph{aṟasṟṟŭ ceyyuka}.) To seize? (\textbf{\ml{പിടിച്ചെടുക്കുക}}, \emph{piṭicceṭukkuka}.) A clue? (\textbf{\ml{സൂചന}}, \emph{sūcana}.) A mark? (\textbf{\ml{അടയാളം}}, \emph{aṭayāḷaṁ}.) A footprint? (\textbf{\ml{കാൽപ്പാട്}}, \emph{kālppāṭŭ}.) A fingerprint? (\textbf{\ml{വിരലടയാളം}}, \emph{viralaṭayāḷaṁ}.) Mysterious? (\textbf{\ml{ദുരൂഹമായ}}, \emph{durūhamāya}.) A crime? (\textbf{\ml{കുറ്റകൃത്യം}}, \emph{kuṟṟakṛtyaṁ}.) A confession? (\textbf{\ml{കുറ്റസമ്മതം}}, \emph{kuṟṟasammataṁ}.) To capture? (\textbf{\ml{പിടികൂടുക}}, \emph{piṭikūṭuka}.) To kidnap? (\textbf{\ml{തട്ടിക്കൊണ്ടുപോകുക}}, \emph{taṭṭikkoṇṭupōkuka}.) To keep watch on? (\textbf{\ml{നിരീക്ഷിക്കുക}}, \emph{nirīkṣikkuka}.) A court of law? (\textbf{\ml{കോടതി}}, \emph{kōṭati}.) A court case? (\textbf{\ml{കേസ്}}, \emph{kēsŭ}.) A trial? (\textbf{\ml{വിചാരണ}}, \emph{vicāraṇa}.) A judge? (\textbf{\ml{ന്യായാധിപൻ}}, \emph{nyāyādhipan}.) The accused? (\textbf{\ml{പ്രതി}}, \emph{prati}.) A verdict? (\textbf{\ml{വിധി}}, \emph{vidhi}.) The guilty one? (\textbf{\ml{കുറ്റക്കാരൻ}}, \emph{kuṟṟakkāran}.) Bail? (\textbf{\ml{ജാമ്യം}}, \emph{jāmyaṁ}.) An innocent person? (\textbf{\ml{നിരപരാധി}}, \emph{niraparādhi}.) To pass sentence? (\textbf{\ml{വിധിക്കുക}}, \emph{vidhikkuka}.) Testimony? (\textbf{\ml{മൊഴി}}, \emph{moḻi}.) An eyewitness? (\textbf{\ml{ദൃക്സാക്ഷി}}, \emph{dṛksākṣi}.) A lie? (\textbf{\ml{കള്ളം}}, \emph{kaḷḷaṁ}.) To testify? (\textbf{\ml{മൊഴി} \ml{നൽകുക}}, \emph{moḻi nalkuka}.) To swear an oath? (\textbf{\ml{സത്യം} \ml{ചെയ്യുക}}, \emph{satyaṁ ceyyuka}.) A jail? (\textbf{\ml{ജയിൽ}}, \emph{jayil}.) Imprisonment? (\textbf{\ml{തടവ്}}, \emph{taṭavŭ}.) A prisoner? (\textbf{\ml{തടവുകാരൻ}}, \emph{taṭavukāran}.) To release? (\textbf{\ml{വിട്ടയയ്ക്കുക}}, \emph{viṭṭayaykkuka}.) To free? (\textbf{\ml{മോചിപ്പിക്കുക}}, \emph{mōcippikkuka}.) A fraud? (\textbf{\ml{തട്ടിപ്പ്}}, \emph{taṭṭippŭ}.) A swindler? (\textbf{\ml{തട്ടിപ്പുകാരൻ}}, \emph{taṭṭippukāran}.) To deceive? (\textbf{\ml{വഞ്ചിക്കുക}}, \emph{vañcikkuka}.) To trick? (\textbf{\ml{പറ്റിക്കുക}}, \emph{paṟṟikkuka}.) To fool? (\textbf{\ml{കബളിപ്പിക്കുക}}, \emph{kabaḷippikkuka}.)
\end{tcolorbox}
